%------------------------------------------------------------------------------%
\section{Aproximações por Polinômios de Taylor}

Uma das principais aplicações das Séries de Taylor é calcular aproximações
para uma função com o auxilio dos Polinômios de Taylor. Se a série converge
temos a fórmula para o erro do Teorema de Taylor~\ref{teo:teorema-taylor}
para estimarmos o erro cometido ou determinarmos o número de termos necessários
para garantir que o erro é menor do que a tolerância aceitável.
Apresentamos a seguir alguns exemplos desse uso das Séries de Taylor.

%------------------------------------------------------------------------------%
\begin{example}
	Escreva \(\int_0^ye^{-x^2}dx\) como uma série de potências centrada em $x=0$.
	Use a série obtida para encontrar uma aproximação para
	\(\int_0^1e^{-x^2}dx\) somando os quatro primeiros termos da série.

	\solution
	Primeiro vamos encontrar a série de Taylor de $f(x)=e^x$.
	Como a $k$-ésima derivada $f^{(k)}(x)=e^x$ e $e^0=1$, segue que
	\[
    e^x =\sum_{n=0}^k\frac{x^n}{n!}+ R_{k}(x)   \;.
  \]
	Verificamos que $R_k(x)\rightarrow 0$ para todo $x\in \R$, pois $f^{(k)}(c)=e^c$ para todo
	$k$, ou seja, as derivadas são limitadas por $\abs{e^c}$.
	Então, para todo $x\in \R$, temos que
	\[
    e^x=\sum_{n=0}^\infty\frac{x^n}{n!}   \;.
  \]
	Seja $u(x)=-x^2$. Note que $u\colon\R\rightarrow\R$ é uma função contínua,
  então, para todo $x\in\R$ temos que
	\[
    e^{-x^2}=\sum_{n=0}^\infty\frac{(-x^2)^n}{n!}=\sum_{n=0}^\infty\frac{(-1)^nx^{2n}}{n!}   \;.
  \]
	Dessa forma,
	\begin{align*}
	  \int_0^ye^{-x^2}dx
    & =\sum_{n=0}^\infty\frac{(-1)^n}{n!}\int_0^y x^{2n} dx    \\[3mm]
	  & =\sum_{n=0}^\infty\frac{(-1)^n}{n!}\frac{y^{2n+1}}{2n+1} \;.
	\end{align*}
	Usando a série obtida, temos
	\begin{align*}
	  \int_0^1e^{-x^2}dx
    & =\sum_{n=0}^\infty\frac{(-1)^n}{n!(2n+1)}             \\[3mm]
	  & \approx \sum_{n=0}^3\frac{(-1)^n}{n!(2n+1)}           \\[3mm]
	  & = 1-\frac{1}{3}+\frac{1}{2!\,5}-\frac{1}{3!\,7}       \\[3mm]
	  & = \frac{26}{35}                                       \\[3mm]
    & \approx 0,7429 \;.
	\end{align*}
\end{example}
%------------------------------------------------------------------------------%

No último exemplo encontramos uma aproximação para o valor de
\[
  \int_0^1e^{-x^2}dx
\]
entretanto não sabemos o erro cometido ao fazer essa aproximação.
Como se trata de uma série alternada, podemos usar a
Proposição~\ref{pro:erro-serie-alternada}
que nos dá uma estimativa superior para esse erro.

%------------------------------------------------------------------------------%
\begin{example}
  Estimar o erro cometido na aproximação do exemplo anterior.
  \solution
	Vimos no último exemplo que
	\[
  	\int_0^1e^{-x^2}dx \approx 0.7429   \;.
	\]
	Com o uso da Proposição~\ref{pro:erro-serie-alternada}
  temos o seguinte limitante para o erro
	\[
  	R_3 \leq a_4
    = \frac{(-1)^4}{4!(2\cdot 4+1)}
    = \frac{1}{4!\,9}
    \approx 0,0046   \;.
	\]
\end{example}
%------------------------------------------------------------------------------%

%------------------------------------------------------------------------------%
\begin{example}
  Aproxime o valor de
  \(
    \sqrt[3]{\sfrac{3}{2}\,}
  \)
  usando Polinômio de Taylor de grau 3 da função
  \(
    f(x)=\sqrt[3]{1+x}
  \).
  Obtenha uma estimativa para o erro cometido.
  \solution
  Podemos construir a Série de Taylor para a função \(f(x)=\sqrt[3]{1+x}\)
  utilizando a Série Binomial
	\[
    \sqrt[3]{1+x} = (1+x)^{\sfrac{1}{3}}
                  = \sum_{n=0}^\infty \binom{\sfrac{1}{3}}{n} x^n
  \]
	onde
  \begin{align*}
    \binom{\sfrac{1}{3}}{0} & = 1                                                                     \\
    \binom{\sfrac{1}{3}}{1} & = \frac{1}{3}                                                           \\
    \binom{\sfrac{1}{3}}{2} & = \frac{\sfrac{1}{3}(\sfrac{1}{3}-1)}{2} = -\frac{1}{9}                 \\
    \binom{\sfrac{1}{3}}{3} & = \frac{\sfrac{1}{3}(\sfrac{1}{3}-1)(\sfrac{1}{3}-2)}{6} = \frac{5}{81} \;.
  \end{align*}
	Assim
	\[
  	\sqrt[3]{1+x} \approx P_3(x) = 1 + \frac{x}{3} - \frac{x^2}{9} + \frac{5x^3}{81} \;.
	\]
	Tomando $x=\sfrac{1}{2}$, temos
	\[
  	\sqrt[3]{\frac{3}{2}\,} = \sqrt[3]{1+\frac{1}{2}}
                           	\approx 1 + \frac{1}{6} - \frac{1}{36} + \frac{5}{648}
                          	\approx 1,1466 \;.
	\]
  Como a série é alternada podemos estimar o erro cometido por
	\[
  	R_3 \leq a_4
      = \abs{\binom{\sfrac{1}{3}}{4}} \; \left(\frac{1}{2}\right)^4
        \approx 0,0026 \;.
	\]
\end{example}
%------------------------------------------------------------------------------%

\begin{figure}
  \begin{center}%
    \vspace{-0.2\baselineskip}%
    \scalebox{0.9}{% Title: gl2ps_renderer figure
% Creator: GL2PS 1.4.2, (C) 1999-2020 C. Geuzaine
% For: Octave
% CreationDate: Tue Sep 28 14:08:47 2021
\setlength{\unitlength}{1pt}
\begin{picture}(0,0)
\includegraphics[width=17cm,height=7cm]{fig_taylor_cos}
\end{picture}%
\begin{picture}(481,198)(0,0)
\fontsize{10}{0}\selectfont\put(19.2756,95.2752){\makebox(0,0)[t]{\textcolor[rgb]{0.15,0.15,0.15}{{$-3 \pi$}}}}
\fontsize{10}{0}\selectfont\put(91.5591,95.2752){\makebox(0,0)[t]{\textcolor[rgb]{0.15,0.15,0.15}{{$-2 \pi$}}}}
\fontsize{10}{0}\selectfont\put(163.843,95.2752){\makebox(0,0)[t]{\textcolor[rgb]{0.15,0.15,0.15}{{$-\pi$}}}}
\fontsize{10}{0}\selectfont\put(308.409,95.2752){\makebox(0,0)[t]{\textcolor[rgb]{0.15,0.15,0.15}{{$\pi$}}}}
\fontsize{10}{0}\selectfont\put(380.693,95.2752){\makebox(0,0)[t]{\textcolor[rgb]{0.15,0.15,0.15}{{$2 \pi$}}}}
\fontsize{10}{0}\selectfont\put(452.976,95.2752){\makebox(0,0)[t]{\textcolor[rgb]{0.15,0.15,0.15}{{$3 \pi$}}}}
\fontsize{10}{0}\selectfont\put(231.13,17.9742){\makebox(0,0)[r]{\textcolor[rgb]{0.15,0.15,0.15}{{-1}}}}
\fontsize{10}{0}\selectfont\put(231.13,187.538){\makebox(0,0)[r]{\textcolor[rgb]{0.15,0.15,0.15}{{1}}}}
\end{picture}
}%
    \vspace{-\baselineskip}%
  \end{center}
  \caption{Comparação da raiz cúbica com seu polinômio de Taylor.}
  \label{fig:funcao_tayplor_sqrt3}
\end{figure}

A Figura~\ref{fig:funcao_tayplor_sqrt3} mostra o gráfico da função
\(f(x)=\sqrt[3]{1+x}\) e o Polinômio de Taylor $P_3$,
assim como a aproximação de $f(\sfrac{1}{2})$ por $P_3(\sfrac{1}{2})$.

%------------------------------------------------------------------------------%
